\section{Introduction}\label{sec:introduction}
A \emph{Littlewood polynomial of length $N$} is a polynomial
\[
 P(z)=\sum_{k=0}^{N-1}\varepsilon_kz^k,
 \qquad \varepsilon_k\in\{-1,1\}.
\]
Write $\T=\R/\Z$ and $\e(t)=\exp(2\pi i t)$, and give $\T$ probability
Haar measure. Parseval's identity gives
\[
 \int_{\T}|P(\e(t))|^2\,dt=N,
 \qquad
 \norm{P}_\infty:=\max_{t\in\T}|P(\e(t))|\ge\sqrt N.
\]
We ask whether a maximum arbitrarily close to this lower bound can
coexist with a minimum bounded below by a fixed positive multiple of
$\sqrt N$.

\begin{theorem}\label{thm:main}
For every $\eta>0$ there is an integer $N_0=N_0(\eta)$ such that, for every
integer $N\ge N_0$, there exist real signs
$\varepsilon_0,\ldots,\varepsilon_{N-1}$ for which
\[
 \frac1{16}\sqrt N
 \le\left|\sum_{k=0}^{N-1}\varepsilon_kz^k\right|
 \le(1+\eta)\sqrt N
 \qquad (|z|=1).
\]
\end{theorem}

The lower constant is independent of the accuracy of the upper bound.
The conclusion holds for a single polynomial at every sufficiently large
length, on the entire circle, including the real points $1$ and $-1$.

\subsection{Context and antecedents}
Questions about the flatness of polynomials with restricted coefficients
have a long history in harmonic analysis. The two-sided flatness question
appears in Erd\H{o}s's problems \cite[Problem 26]{Erdos1957} and in
Littlewood's work \cite{Littlewood1966}. The Rudin--Shapiro construction
of real-sign polynomials gives an upper bound $\sqrt{2N}$ at dyadic lengths
\cite{Rudin1959}. The requirement that the coefficients be real signs is
substantial: Kahane constructed polynomials with complex unimodular
coefficients whose normalized modulus tends uniformly to one
\cite{Kahane1980}; see also Bombieri and Bourgain
\cite{BombieriBourgain2009}. The real-sign questions and their relation to
the Parseval lower bound are discussed in
\cite[Problems 4.13 and 4.31]{HaymanLingham2018}.

Balister, Bollob\'as, Morris, Sahasrabudhe and Tiba proved that flat
Littlewood polynomials exist at every degree at least two: their modulus lies between
two absolute positive multiples of the square root of the length
\cite[Theorem 1.1]{BalisterEtAl2020}. The companion manuscript
\cite{OpenAIMinimal} establishes that the best possible upper constant
tends to one through all lengths. Theorem~\ref{thm:main} combines the
nearly minimal upper bound with a fixed positive lower bound. These two
properties must be achieved by the same choice of signs, so neither of the
preceding existence statements supplies the conclusion by itself.

The construction in \cite{OpenAIMinimal} provides more than an upper bound.
It gives relaxed real coefficients in $[-1,1]$ whose Fourier polynomial
is uniformly approximated by quadratic-phase waves with disjoint supports.
The modulus of this leading term has almost full mean-square mass, but it vanishes
on some arcs. We reproduce the auxiliary polynomial, interval packing,
sampling, and rounding arguments needed here. The standard external
inputs are the Pippenger--Spencer hypergraph-coloring theorem and the
Lovett--Meka partial-coloring theorem, stated at their points of use.

Our additional step constructs a correction on the arcs where the leading
amplitude is small. On each arc its phase is $N$ times a fixed piecewise
quadratic function, plus an affine adjustment with slope bounded
independently of $N$. We call the fixed function the leading phase.
At a transition, the correction and existing wave have the same leading
phase derivative, and the affine adjustment makes their full phase values
agree modulo integers. This keeps the waves aligned while the correction turns on.
Across most of each arc, the leading phase derivative runs through an
interval assigned only to that arc and its reflection. This separation controls individual
Fourier coefficients; keeping that derivative strictly between zero and
one controls the Fourier mass
outside the permitted degrees. The two estimates preserve both the
coefficient constraint and a lower-modulus margin.

The proof of \cite{BalisterEtAl2020} also repairs small-modulus regions:
a cosine polynomial from the Rudin--Shapiro construction is supplemented
by a sine polynomial that is large on the exceptional arcs, with
discrepancy estimates controlling the coefficients. Here the correction
acts on the quadratic waves of \cite{OpenAIMinimal}, and its coefficient
cost tends to zero as the upper bound approaches its minimum.

\subsection{Proof strategy and notation}
For a real vector $X=(X_0,\ldots,X_{N-1})$, write
\[
 U_X(t)=\frac1{\sqrt N}\sum_{k=0}^{N-1}X_k\e(kt),
 \qquad
 \mu(X)=\frac12\sum_{k=0}^{N-1}(1-|X_k|)
 \quad\text{if }X\in[-1,1]^N.
\]
The quantity $\mu(X)$ measures how far the coefficients are from signs.
The rounding estimate proved in Section~\ref{sec:rounding} changes $U_X$
by at most
\[
 \frac C{\sqrt N}
 +C\sqrt{\frac{\mu(X)}N
          \log\!\left(\frac{80N}{\mu(X)}\right)},
\]
with the square-root term defined to be zero at $\mu(X)=0$.
Thus a small relative defect permits a uniformly small change on the whole
circle. The construction must create its positive lower bound before this
rounding is applied.

Fix a small accuracy parameter $\delta>0$. Sections~\ref{sec:spreading}
and~\ref{sec:packing} construct a bounded real trigonometric polynomial
on an auxiliary torus and pack the angle intervals that will support
the resulting quadratic waves. Section~\ref{sec:sampling} produces vectors $Y_N\in[-1,1]^N$ and
leading functions $B_N$ with
\[
 U_{Y_N}=B_N+o(1),\qquad |B_N|=B\le1+O(\delta),\qquad
 \int_{\T}B^2\ge1-O(\delta).
\]
Here $B$ is a fixed sum of disjoint smooth bumps, independent of
$N$. In particular, the arcs where $B<1/8$ have total length $O(\delta)$.

Section~\ref{sec:correction} constructs a conjugate-symmetric function
$R_N$ on those arcs such that $B_N+R_N$ has minimum at least $1/8-o(1)$,
without raising its maximum. The essential Fourier estimates are
\[
 \sqrt N\,|\widehat R_N(k)|=O(\sqrt\delta)\quad(0\le k<N),
 \qquad
 \sum_{k\notin\{0,\ldots,N-1\}}|\widehat R_N(k)|=o(1).
\]
The first bound permits a small rescaling that returns the corrected real
coefficients to $[-1,1]$. The second shows that restricting the correction
to the required degrees preserves its uniform effect on the circle.
The rescaled coefficients still have relative defect tending to zero with
$\delta$. Section~\ref{sec:completion} applies the rounding estimate and
chooses the parameters to prove Theorem~\ref{thm:main}.

All torus integrals below use probability Haar measure. For an integrable
function $f$ on $\T^m$ we use
\[
 \widehat f(a)=\int_{\T^m}f(y)\e(-a\cdot y)\,dy
 \qquad(a\in\Z^m).
\]
All auxiliary dimensions, frequencies, intervals, and smooth cutoffs are
fixed after $\delta$ is chosen. Unless another limit is specified, $o(1)$
denotes a quantity tending to zero as $N\to\infty$ with these data fixed.
Constants denoted by $O_\delta$ may depend on these fixed data.
An \emph{absolute} constant is independent of both $\delta$ and $N$.
